\section{Open Challenges and Limitations}\label{sec:challenges}

\subsection{Synthesized Open Research Challenges (\texorpdfstring{\texttt{OC1--OC7}}{OC1-OC7})}\label{subsec:open_challenges}
Combining the literature gap matrix (Section~\ref{sec:taxonomy}) with the simulation trade-offs in Section~\ref{sec:results} identifies seven open research challenges:
\begin{itemize}
    \item \textbf{OC1 (Cross-Region KV-Cache Transport):} Even at 100~Gbps WAN, the carbon-and-latency gate prunes cross-region P/D transfers for over 99.9\% of arrivals. Closing that gap requires semantic token-level KV compression~\cite{jiang2026kvcache} integrated with live grid-carbon price gates.
    \item \textbf{OC2 (Live Marginal vs.\ Average Carbon Control):} Most cloud systems, including the baseline \texttt{E1} sweeps, optimize average regional grid intensity ($\text{gCO}_2\text{eq/kWh}$). Dispatching multi-megawatt LLM bursts requires marginal emission factors and grid-curtailment telemetry to reflect fossil peaker dispatch.
    \item \textbf{OC3 (Sub-Minute Embodied Cluster Sizing):} While per-request dispatch in \texttt{P4} minimizes operational carbon plus queue risk, \texttt{H4} shows that idle over-provisioning in clean zones reverses lifecycle gains (+32.6\% $\text{SCI}_{\text{tot}}$). Joint control planes must couple sub-second routing with hourly embodied scale-down and heterogeneous reuse (9~yr/3~yr lifetimes~\cite{li2025ecoserve,hewage2025aging}).
    \item \textbf{OC4 (Serverless Multi-Tenant Checkpoint Sharing):} On \texttt{TEPCO}, dynamic keep-alive still faces an 8.0\% vs.\ 28.60\% carbon-SLO squeeze. Page-shared GPU weight caching across concurrent LoRA adapters and MPS partitions~\cite{wang2026festina} can amortize 10.3~s NVMe cold starts across tenants.
    \item \textbf{OC5 (Water-Energy-Carbon Co-Optimization):} As CASA~\cite{qi2024casa} notes for serverless CPUs, shifting LLM inference purely toward dry wind/solar regions risks shifting cooling water stress unless water intensity enters Eq.~\eqref{eq:sci}.
    \item \textbf{OC6 (Bursty Code-Completion Admission Control):} Under $\text{CV} = 13.152$, no fixed 6-node strategy stayed below 5\% SLO violations without elastic spillover or prefill-priority admission pricing.
    \item \textbf{OC7 (Hardware Testbed Verification Across Continents):} Only 28.6\% of reviewed empirical systems pair testbeds with trace simulators, and none tests live WAN disaggregated LLM serving across continents (0/6); multi-continent hardware validation remains untested.
\end{itemize}

\subsection{Limitations and Threats to Validity}\label{subsec:limitations}
Several methodological constraints bound the interpretation of these results:
\begin{enumerate}
    \item \textbf{Simulation Abstraction vs.\ Physical Clusters:} Findings rely on analytical queueing and SimPy 4.1.2 discrete-event execution; no measurements come from instrumented multi-continent Kubernetes/vLLM deployments. Actual RDMA tail jitter, PCIe thermal throttling, and Linux memory fragmentation may shift absolute latency margins, although relative policy rankings (\texttt{P4} over \texttt{P3} and \texttt{B1}/\texttt{B2}) rest on structural bottlenecks observed in physical single-cluster studies~\cite{patel2024splitwise,zhong2024distserve,li2026revisiting}.
    \item \textbf{Calibrated Model Parameters:} Cross-region WAN RTT (20--200~ms), commodity WAN bandwidth (10--100~Gbps), sublinear MFU exponent ($\gamma=0.70$), and cold-start GPU power fraction ($f_{\text{cold}}=0.65$) are engineering assumptions bounded by published ranges (Table~\ref{tab:calibrations}) and not direct hardware measurements; \texttt{SW3} and \texttt{SW4} sweep these parameters across wide intervals.
    \item \textbf{24-Hour Trace Tiling:} Raw Azure LLM traces span 58.36 min (\texttt{conv}) and 57.27 min (\texttt{code})~\cite{azure2024trace}. Multi-day evaluations rely on block resampling across 24-hour UK/Japan grid curves while holding exact token distributions and arrival CVs constant.
    \item \textbf{Live Dispatch vs.\ Provisioning Scope:} Inside Eq.~\eqref{eq:dispatch}, per-request \texttt{P4} routing scores operational cost and queue pressure; hardware embodied amortization is evaluated across fixed versus over-provisioned cluster configurations (\texttt{H4\_elastic}) without dynamically powering servers on or off inside \texttt{E1}.
    \item \textbf{Corpus Access and Preprints:} Five works (10.6\%) could be verified only via abstracts~\cite{roy2026lowcarb,calavaro2026offloading,gupta2022act,sinha2025utilization,calheiros2011cloudsim}, and eight full-text works (17.0\%) are arXiv or accepted conference preprints (6 arXiv-only, 2 accepted preprints), which may undergo peer-reviewed revision prior to archival publication.
\end{enumerate}
